\chapter{Authorization on the Wire}
\label{ch:authorization}

\epigraph{``You shall not pass!''}{Gandalf, \emph{The Fellowship of the Ring}
(2001)}

Gandalf's line is the most famous authorization decision in cinema, and it is
worth noticing what it actually was. He did not check a credential. He asserted
a boundary, at a chokepoint, against something that had every intention of
walking straight through. That is the job. The credential is just the paperwork.

This is a security chapter with a performance chapter hiding inside it, which is
the same structure \emph{High Performance Browser Networking} used for TLS, and
for the same reason. Auth is not a policy you configure once. It is a sequence
of network round trips that happen before your application does any work at all,
and if you do not amortise them you will pay for them forever.

\section{The shape of it}

MCP does not invent an auth scheme. It profiles OAuth 2.1 and a stack of RFCs
around it, taking a deliberately small subset. The roles map cleanly:

\begin{description}[leftmargin=1.4em, style=nextline, itemsep=3pt]
\item[The MCP server is an OAuth resource server.] It accepts access tokens and
validates them. It does not issue them.
\item[The MCP client is an OAuth client.] It obtains tokens on behalf of a user.
\item[The authorization server is somebody else's problem.] Entra ID, Okta,
Keycloak, or your own. MCP specifies how a client \emph{finds} it and says nothing about its
internals.
\end{description}

Two scoping rules before we go further.

Authorization is \textbf{optional}. Plenty of MCP servers need none.

Authorization is \textbf{for HTTP}. On stdio the specification says explicitly
that you should \emph{not} do this, and should take credentials from the
environment instead. The client launched the process; it can hand it an API key
through the environment block. Running an OAuth flow against your own subprocess
is theatre.

\begin{py}
{
  "mcpServers": {
    "meridian-risk": {
      "command": "python3",
      "args": ["-m", "meridian.serve", "risk"],
      "env": { "MERIDIAN_API_KEY": "..." }
    }
  }
}
\end{py}

\section{The full flow, and its price tag}

\bookfig[1.0]{oauth-flow}{OAuth 2.1 as MCP profiles it. Count the round trips on
the left. Every one of them happens before your server does a single unit of
useful work.}{oauth-flow}

Walking it:

\textbf{1. The challenge.} The client makes an unauthenticated request. The
server answers 401 with a \texttt{WWW-Authenticate} header that points at its
metadata and, ideally, names the scopes needed:

\begin{http}
HTTP/1.1 401 Unauthorized
WWW-Authenticate: Bearer resource_metadata="https://mcp.meridian.example/.well-known/oauth-protected-resource",
                         scope="risk:read"
\end{http}

That \texttt{scope} parameter is a courtesy with real value. Without it the
client falls back to requesting everything in \texttt{scopes\_supported}, which
means the user sees a consent screen asking for more than the operation needs,
which means the user learns to click through consent screens.

\textbf{2. Protected resource metadata.} Required, by RFC 9728. The client
fetches it and learns which authorization servers this resource trusts.

\textbf{3. Authorization server metadata.} The client tries OAuth 2.0
Authorization Server Metadata and OpenID Connect Discovery in priority order.
Clients must support both, because the ecosystem is split and will stay split.

\textbf{4. Authorization.} The client opens a browser at the authorization
server's authorization endpoint. The user signs in, sees a consent screen, and
approves. The browser is then redirected to a URI the client is listening on,
carrying an authorization code. Two parameters on that request are mandatory,
and both are worth understanding rather than copying.

The first is \texttt{code\_challenge}.

\textbf{5. Token exchange.} The client posts the code back, together with the
PKCE verifier and the \texttt{resource} parameter, and receives an access token.

Then, finally, the original request runs.

\subsection{PKCE, and why an authorization code is not enough}

PKCE (pronounced ``pixie'', Proof Key for Code Exchange, RFC 7636) exists
because an authorization code travels through a browser redirect, and a browser
redirect is not a private channel.

Consider where the code lands. A desktop or CLI client cannot receive a redirect
at a public HTTPS address, so it starts a tiny local web server and registers
\texttt{http://127.0.0.1:3000/callback} as its redirect URI. Any other process
on that machine can bind that port first and catch the redirect. On mobile the
equivalent trick is registering the same custom URI scheme. Either way an
attacker who catches the code can walk up to the token endpoint and redeem it,
and the user's token comes back to somebody else.

PKCE binds the code to whoever asked for it, in three steps.

\begin{enumerate}[leftmargin=1.6em, itemsep=3pt]
\item Before redirecting, the client generates a \texttt{code\_verifier}: a
high-entropy random string of 43 to 128 characters. It keeps this in memory,
associated with this one authorization attempt, and never transmits it during
the redirect.
\item It sends the SHA-256 hash of that verifier, base64url-encoded, as
\texttt{code\_challenge}, along with \texttt{code\_challenge\_method=S256}. The
authorization server files the challenge next to the code it issues.
\item At the token endpoint the client presents the original verifier. The
server hashes it and compares it to the stored challenge. A mismatch means the
party redeeming the code is not the party that requested it, and the exchange
is refused.
\end{enumerate}

\begin{plain}
code_verifier   dBjftJeZ4CVP-mB92K27uhbUJU1p1r_wW1gFWFOEjXk
code_challenge  E9Melhoa2OwvFrEMTJguCHaoeK1t8URWbuGJSstw-cM
                = base64url(sha256(code_verifier))
\end{plain}

An intercepted code is now inert. The interceptor holds the code and does not
hold the verifier, and the verifier never left the client's memory. The hash is
one-way, so the challenge that did travel over the redirect gives nothing away.

One requirement here is easy to skip, and the specification states it as a MUST:
the client must \emph{verify} that the authorization server supports PKCE before
it begins, by looking for \texttt{code\_challenge\_methods\_supported} in the
metadata it fetched in step 3. If that field is absent, the client must refuse
to proceed. Sending a challenge to a server that ignores challenges produces a
flow that looks identical, feels identical, and protects nothing, which is a
worse outcome than a visible failure. OpenID Connect Discovery does not
technically define that field, and MCP requires it anyway; an OpenID provider
that omits it is unusable for MCP.

\subsection{One bookkeeping step with no network cost}

While the browser is open, the client records the \texttt{issuer} value from the
authorization server metadata it validated in step 3, and stores it in the same
per-request record that holds the PKCE verifier and the \texttt{state} value.

Nothing uses it yet. \secref{issuer-validation} is where it earns its keep, and
the requirement that it come from \emph{validated} metadata is the whole point:
an expected issuer read from an untrusted source proves nothing when you compare
against it later.

\begin{measurebox}{What the handshake costs (modelled)}
\begin{tabular}{@{}lrr@{}}
\toprule
\thead{Step} & \thead{Round trips} & \thead{Typical} \\
\midrule
401 challenge                        & 1 & 40\,ms \\
Protected resource metadata          & 1 & 40\,ms \\
Authorization server metadata        & 1 & 60\,ms \\
Authorization (includes a human)     & 1 & seconds \\
Token exchange                       & 1 & 80\,ms \\
\midrule
\textbf{Machine time before first work} & \textbf{4} & \textbf{\textasciitilde{}220\,ms} \\
\bottomrule
\end{tabular}

\vspace{6pt}
\footnotesize
Modelled at 40\,ms per hop. Meridian's benchmark auth runs locally and would
flatter these badly, so these are estimates. The number that matters is the round-trip count,
which is exact.
\end{measurebox}

Four machine round trips, plus a human. If you pay that per request you have
built something unusable. If you pay it once per session and cache aggressively,
it disappears. That is the entire performance content of this chapter, and it is
worth stating as a rule:

\keyidea{Steps 1 through 3 are properties of the server, so cache them per
origin, for hours. The token is a property of the user and the scope set, so
cache it until it expires. Refresh it while it still works.}

\section{Client identity in 2026}

Here is the part that changed, and it changed for a good reason.

MCP's awkward case is a client and a server with no prior relationship. A user
installs some MCP client, points it at some MCP server, and expects it to work.
There is no administrator to pre-register an OAuth client.

The 2025 answer was Dynamic Client Registration: the client POSTs to a
registration endpoint and gets credentials. It works, and it has a nasty
property. Every client instance registers separately, so an authorization server
accumulates registrations forever, most of them belonging to a laptop that was
reimaged eight months ago. It is a write endpoint open to the internet that
grows without bound.

\subsection{Client ID Metadata Documents}
\label{sec:client-identity}

The replacement is neat. The client id \emph{is} an HTTPS URL, and that URL
serves the client's metadata:

\begin{wire}
{
  "client_id": "https://app.example.com/oauth/client-metadata.json",
  "client_name": "Example MCP Client",
  "client_uri": "https://app.example.com",
  "logo_uri": "https://app.example.com/logo.png",
  "redirect_uris": [
    "http://127.0.0.1:3000/callback",
    "http://localhost:3000/callback"
  ],
  "grant_types": ["authorization_code"],
  "response_types": ["code"],
  "token_endpoint_auth_method": "none"
}
\end{wire}

The authorization server sees a URL-shaped \texttt{client\_id}, fetches it,
validates that the document's \texttt{client\_id} matches the URL exactly,
checks the \texttt{redirect\_uri} against the list, and proceeds.

What this buys:

\begin{itemize}[leftmargin=1.6em, itemsep=3pt]
\item \textbf{No registration round trip.} The zero-numbered step in
\figref{oauth-flow} happens on the authorization server's own time, and it is
cacheable under ordinary HTTP cache headers.
\item \textbf{No unbounded registration table.} One document serves every
instance of the client.
\item \textbf{Portability.} The same client id works at every authorization
server, because it resolves on demand. Credentials from DCR are bound to the
issuer that minted them and must be re-registered when it changes.
\end{itemize}

Servers advertise support with one metadata field:

\begin{wire}
{ "client_id_metadata_document_supported": true }
\end{wire}

The client's selection order, when it supports everything:

\begin{enumerate}[leftmargin=1.6em, itemsep=2pt]
\item Pre-registered credentials, if it has them for this server
\item Client ID Metadata Documents, if the AS advertises support
\item Dynamic Client Registration, if the AS has a \texttt{registration\_endpoint}
\item Ask the user to paste a client id
\end{enumerate}

\begin{legacybox}{Dynamic Client Registration (RFC 7591)}
Deprecated as of \texttt{2026-07-28}, eligible for removal no earlier than
2027-07-28. Retained because plenty of 2025-era authorization servers speak
nothing else.

If you must use it, specify \texttt{application\_type} explicitly. Omitting it
defaults to \texttt{"web"} under OpenID Connect, which rejects the
\texttt{localhost} redirect URIs that desktop and CLI clients need. Use
\texttt{"native"} for anything local. This is a genuinely common failure and the
error message is never helpful.

Also: key persisted credentials by the authorization server's \texttt{issuer}.
Reusing a client id across authorization servers is a specification violation,
and it fails in confusing ways.
\end{legacybox}

\section{Two validations that stop real attacks}

Most of OAuth's complexity is scar tissue. Two pieces of it are worth understanding in detail, because both fail silently
when you get them wrong.

\subsection{Audience binding, or: do not be a confused deputy}
\label{sec:audience-binding}

An access token carries an \emph{audience}: the identity of the resource server
the token was minted for. In a JWT it is the \texttt{aud} claim; with opaque
tokens it comes back from introspection. It answers a question distinct from
``is this token real'', and the distinction is the whole subsection:

\begin{plain}
iss  who issued this token          https://login.example.com
sub  which user it speaks for       user-8891
aud  which server it is for         https://mcp.meridian.example
exp  when it stops working          1793404800
\end{plain}

A token can be perfectly valid, unexpired, correctly signed by an issuer you
trust, and speak for a real user, and still not be \emph{for you}. Checking the
signature answers the first four questions. Only checking \texttt{aud} answers
the fifth.

Audience binding is the mechanism that puts the right value in that claim. It is
RFC 8707, and it has two halves.

The client half: send a \texttt{resource} parameter naming the server the token
is for, on both the authorization request and the token request.

\begin{plain}
&resource=https%3A%2F%2Fmcp.meridian.example
\end{plain}

The authorization server copies that into the token's audience. The server half:
validate, on every request, that the token presented was issued for you. A token
whose audience names somebody else gets a 401, no matter how well-formed it is.

Skip the server half and you have built a confused deputy.

\begin{dangerbox}{The confused deputy}
A confused deputy is a program with more authority than its caller, tricked into
using that authority on the caller's behalf. The name is from a 1988 paper by
Norm Hardy about a compiler that could write to a billing file its users could
not, and could be talked into overwriting it.

Here is the MCP version. A user installs a malicious MCP server and authorizes
it, willingly, because it claims to summarise their calendar. The authorization
server issues it a token for the user. That server now presents the same token
to \emph{your} server. If you check only that the token is valid, unexpired, and
signed by an issuer you trust, every one of those checks passes, and you go and
move the user's money at the attacker's direction.

The token was real. The user did consent. They consented to something else, and
the audience claim is where that distinction is written down.
\end{dangerbox}

Three rules follow, and the specification states all three as MUST:

\begin{itemize}[leftmargin=1.6em, itemsep=2pt]
\item Clients must not send a token to any server other than the one its
authorization server issued it for.
\item Servers must accept only tokens valid for their own resources.
\item Servers must not accept or forward any other token.
\end{itemize}

That third one forbids token passthrough, and it is worth being blunt: an MCP
server that takes the client's token and replays it at a third-party API is
doing the forbidden thing. Chapter~\ref{ch:mrtr} covers the correct pattern,
which is URL-mode elicitation, and Chapter~\ref{ch:access-control} covers why
the shortcut is so tempting.

\subsection{Issuer validation, RFC 9207}
\label{sec:issuer-validation}

New in this revision, and small enough to skip by accident.

Audience binding protects the server from a token issued for somebody else.
Issuer validation protects the \emph{client} from a code issued by somebody
else, and the attack it stops is called a mix-up.

A mix-up needs one precondition: the client talks to more than one authorization
server, and the attacker controls one of them. That is not exotic. A client that
connects to several MCP servers accumulates several issuers, and any one of them
could be hostile.

The attack runs like this. The user asks the client to connect to the attacker's
MCP server, which is an ordinary thing to do and not yet a compromise. The
client discovers the attacker's authorization server and starts a flow with it.
The attacker's authorization server, instead of authenticating the user itself,
redirects the browser onward to an honest authorization server the client also
uses, with the client's own \texttt{client\_id} and redirect URI. The user sees
a login page they recognise, from a company they trust, and signs in. The code
comes back to the client's redirect URI.

The client is now holding a code issued by the honest server while believing it
is talking to the attacker's server. So it does the natural thing: it sends the
code, and the PKCE verifier, to the attacker's token endpoint. The attacker
takes that code and redeems it at the honest server. PKCE did not help, because
the client handed over the verifier itself.

The fix is one string comparison. Before redirecting, the client records the
\texttt{issuer} from the authorization server's validated metadata, which is the
bookkeeping step from earlier in this chapter. When the authorization response
comes back, it should carry an \texttt{iss} parameter naming who actually issued
the code. The client compares the two before redeeming anything. In the attack
above they differ, and the flow stops before the code moves.

\begin{table}[htbp]
\centering
\small
\begin{tabularx}{\textwidth}{@{}llX@{}}
\toprule
\thead{AS advertises \texttt{iss}} & \thead{\texttt{iss} present} & \thead{Client action} \\
\midrule
true  & yes & Compare, by simple string comparison \\
true  & no  & \textbf{Reject the response} \\
false & yes & Compare anyway \\
false & no  & Proceed \\
\bottomrule
\end{tabularx}
\caption{RFC 9207 validation. Row two is the one that exists specifically to
stop mix-up attacks.}
\label{tab:iss-validation}
\end{table}

Two details people get wrong.

\textbf{Simple string comparison.} After URL-decoding, do not normalise. No case
folding on scheme or host, no default-port elision, no trailing-slash
adjustment, no percent-encoding normalisation. Normalising creates equivalence
classes, and equivalence classes are where attackers live.

\textbf{It applies to error responses too.} If \texttt{iss} does not match, do
not display \texttt{error\_description}. That string is attacker-controlled and
you were about to render it in your UI.

\section{Scopes without making the user hate you}

Scope design is where auth stops being mechanics and starts being product
design.

The principle is least privilege, and the mechanism is step-up: request the
minimum, and ask for more when an operation actually needs it. When a call needs
a scope the token lacks:

\begin{http}
HTTP/1.1 403 Forbidden
WWW-Authenticate: Bearer error="insufficient_scope",
                         scope="risk:write",
                         resource_metadata="https://mcp.meridian.example/.well-known/oauth-protected-resource",
                         error_description="Writing a risk override requires risk:write"
\end{http}

Note the 403. A 401 means ``I do not know who you are''; a 403 means ``I know
exactly who you are and you may not do that''. Clients branch on this,
so getting it wrong sends them into a re-authentication loop that cannot succeed.

Two rules for servers, both about not being annoying:

\textbf{Challenge with everything the operation needs, at once.} Returning one
missing scope, then another on the retry, forces a separate consent round trip
per scope. Determine the full set, then emit it together.

\textbf{Be consistent.} Whatever strategy you pick for constructing scope sets,
apply it uniformly. Clients cache and predict; a server that varies its
challenges defeats both.

And one rule for clients, which is the one that gets forgotten:

\keyidea{On a step-up, request the \emph{union} of what you already had and what
was just challenged. Requesting only the new scope silently drops the old ones,
and the next operation fails in a way that looks like a server bug.}

\section{Failure mid-loop}

Now the operationally interesting part, which the specifications do not really
address because it is an application concern.

An agent loop runs for a while. Tokens expire. What happens when step 7 of a
12-step task gets a 401?

The wrong answer is to fail the task. The user's work is half done, the model
holds context that took real money to build, and throwing it away to re-run an
OAuth flow from scratch is the worst of both worlds.

The right answer treats re-auth as a recoverable interruption:

\begin{py}
def call_with_reauth(self, name, arguments, *, max_attempts=2):
    """Refresh once, retry once, then stop.

    The retry limit matters. A server that returns 401 for a reason the
    client cannot fix (a revoked grant, a disabled account) will do so
    forever, and an unbounded retry turns that into an outage plus a
    denial-of-service against your own authorization server.
    """
    for attempt in range(max_attempts):
        try:
            return self.call_tool(name, arguments)
        except Unauthorized:
            if attempt + 1 == max_attempts:
                raise
            if not self.refresh_token():
                raise            # refresh failed; a human is needed
    raise Unauthorized("re-authentication did not help")
\end{py}

Better still, do not get there. Refresh proactively when a token is within some margin of expiry, so the
refresh happens between loop iterations and never in the middle of one.

\begin{table}[htbp]
\centering
\small
\begin{tabularx}{\textwidth}{@{}lX@{}}
\toprule
\thead{Situation} & \thead{What the host should do} \\
\midrule
Access token expired, refresh token valid
  & Refresh silently. The user sees nothing. \\
Refresh token expired or revoked
  & Pause the loop, keep the context, prompt for re-auth, resume. \\
Insufficient scope (403)
  & Step up with the union of scopes, then retry the one operation. \\
Grant revoked entirely
  & Fail cleanly and say so. Do not retry; you are attacking your own AS. \\
\bottomrule
\end{tabularx}
\caption{Auth failure modes and their handling. Only the last one should ever
lose the task.}
\label{tab:auth-failures}
\end{table}

\section{Enterprise identity providers}

Brief, practical, and mostly reassurance.

You almost certainly do not need to write an authorization server. Entra ID,
Okta, Keycloak, and Auth0 are all OAuth 2.1 capable. The MCP-specific
requirements are narrow:

\begin{enumerate}[leftmargin=1.6em, itemsep=3pt]
\item \textbf{Your MCP server publishes protected resource metadata} at
\texttt{/.well-known/\allowbreak oauth-protected-\allowbreak resource}. This is
a static JSON document. It is the only piece of OAuth infrastructure you are
obliged to build.
\item \textbf{Your IdP supports RFC 8707 resource indicators.} Most do now. If
yours does not, audience binding has to be enforced by convention, which is
weaker and worth pushing back on.
\item \textbf{Your IdP supports CIMD, or you fall back to DCR or
pre-registration.} All three are legitimate; the priority order is in
\secref{client-identity}.
\end{enumerate}

What you should not do is fork your MCP server per identity provider. The
provider is a URL in a metadata document. If provider-specific code is leaking
into your request handlers, something has gone wrong upstream of the handlers.

\section{Auth in Meridian}

Meridian's HTTP transport takes an authenticator callable, so the protocol layer
never learns what a token is:

\begin{py}
def _handle_post(self, h) -> None:
    ...
    auth = None
    if self.authenticator is not None:
        auth = self.authenticator(h.headers.get("Authorization"))
        if auth is None:
            h.send_response(401)
            h.send_header(
                "WWW-Authenticate",
                'Bearer resource_metadata='
                f'"http://{self.host}:{self.port}'
                '/.well-known/oauth-protected-resource"')
            h.send_header("Content-Length", "0")
            h.end_headers()
            return
\end{py}

The resulting claims land on the request context, and from there they reach the
two places that need them:

\textbf{Filtering the catalogue.} The specification is precise here, and the
precision matters. A server's tool list must not vary per \emph{connection}, but
it may vary by the \emph{authorization presented on the request}. Those sound
similar and are completely different: credentials are per-request input, so
filtering on them is stateless, whereas filtering on connection identity is not.

\begin{py}
def visible_tools(self, ctx: RequestContext) -> list[Tool]:
    """Override to filter by scope. The set may vary by authorization,
    never by connection."""
    return sorted(self._tools.values(), key=lambda t: t.name)
\end{py}

\textbf{Binding MRTR state to a principal.} This is the one that stops a real
attack. Meridian seals its \texttt{requestState} with the authenticated subject
inside, and rejects state presented by anybody else:

\begin{py}
# sketch: elided from the shipped handler
return input_required(
    input_requests={"approval": elicit_form(...)},
    request_state=SEALER.seal(
        {"accountId": account_id, "exposureUsd": exposure},
        principal=(ctx.auth or {}).get("sub"),
        method=ctx.method,
        params=ctx.params,
    ),
)
\end{py}

Without the principal binding, Bob could replay Alice's half-completed
high-value approval and inherit it. Chapter~\ref{ch:mrtr} covers the sealing
scheme in full, and Meridian tests each binding separately:

\begin{py}
def test_principal_binding_blocks_cross_user_replay(self):
    token = self.sealer.seal({"txn": "T1"}, principal="alice")
    self.assertEqual(self.sealer.open(token, principal="alice"),
                     {"txn": "T1"})
    with self.assertRaises(McpError):
        self.sealer.open(token, principal="bob")
\end{py}

\begin{notebox}{Cache scope and authorization are the same conversation}
A \texttt{private} cached response may only be reused within the same
authorization context. A \texttt{public} one may be served to anyone, including
across tokens, even when it came from an authenticated endpoint.

So \texttt{cacheScope} is an access-control decision wearing performance
clothing. Meridian keys \emph{every} cache entry by principal, public or not,
because the cost of a duplicate entry is a few kilobytes and the cost of being
wrong is a data breach. Chapter~\ref{ch:resources} has the measurements.
\end{notebox}

\section{What to remember}

\begin{itemize}[leftmargin=1.6em, itemsep=3pt]
\item OAuth 2.1 with PKCE, for HTTP only. On stdio, use the environment.
\item Four machine round trips before any work happens. Cache discovery per
origin, cache tokens until expiry, refresh before you need to.
\item Client ID Metadata Documents replace Dynamic Client Registration. The
client id is an HTTPS URL that resolves to metadata.
\item Audience binding via RFC 8707 is what stops the confused deputy. Never
accept or forward a token issued for somebody else.
\item Validate \texttt{iss} before redeeming a code, with simple string
comparison and no normalisation, including on error responses.
\item Challenge with all the scopes an operation needs at once; step up with the
union, never the difference.
\item Treat an expired token as a recoverable interruption. Bound your retries.
\item Tool visibility may vary by authorization, never by connection.
\end{itemize}

That closes Part~I. You now know the physics, the architecture, the wire, and
the credentials. Part~II is what you actually build with them, starting with
the primitive that does the work.
